% SGA 3，第 VIII 篇，英文重建底本说明。
% 范围：续完推论 4.7 的证明，并完成
% 命题 4.6，位于权威源单篇第 14 页顶部
% （合阅本第 630 页）。止于第 5 节标题之前。
% 权威源：Polo--Gille Exp8-8nov09.pdf.
% 权威源单篇第 14 页; 合阅本第 630 页.
这表明，将 \(t\) 送至 \(f\) 就得到分次 \(A_0\)-代数的同构
\[
  A_0[t,t^{-1}]\longrightarrow A
\]
这就完成了 \SGARefLink{sga3:viii:corollary:4:7}{4.7} 的证明。

在 \SGARefLink{sga3:viii:proposition:4:6}{4.6} 的假设下，\SGARefLink{sga3:viii:corollary:4:7}{4.7} 已表明对 \(m\in M\) 的各 \(\mathscr A_m\) 都可逆。因此，为证明条件 \SGARefLink{sga3:viii:proposition:4:1:item:b:01}{4.1(b)}，可以假设 \(\mathscr A_m\) 和 \(\mathscr A_{m'}\) 分别有基 \(f_m\) 和 \(f_{m'}\)，其逆为
\[
  f_m^{-1}\in\Gamma(\mathscr A_{-m})
  \qquad\text{及}\qquad
  f_{m'}^{-1}\in\Gamma(\mathscr A_{-m'}).
\]
于是乘以
\[
  f_m^{-1}f_{m'}^{-1}
  \in\Gamma(\mathscr A_{-m-m'})
\]
定义同态
\[
  \mathscr A_{m+m'}\longrightarrow
  \mathscr A_0\simeq\mathcal O_S
\]
它将 \(f_m\otimes f_{m'}\) 的像送至 \(\mathcal O_S\) 的单位截面。在图
\SGARefTarget{sga3:viii:proposition:4:6:diagram:01}
\[
\begin{tikzcd}[column sep=large]
  \mathscr A_m\otimes\mathscr A_{m'}
    \arrow[r,"u"]
    \arrow[rr,bend left=32,"w"]
  & \mathscr A_{m+m'}
    \arrow[r,"v"]
  & \mathscr A_0\simeq\mathcal O_S ,
\end{tikzcd}
\]
中，映射 \(w\) 和 \(v\) 因而都是可逆层的满态射，故为同构；从而 \(u\) 为同构。
\hfill\(\square\)
